%%%%%%%%%%%%%%%%%%%%%%%%%%%%%%%%%%%%%%%%%
% Masters/Doctoral Thesis 
% LaTeX Template
% Version 2.5 (27/8/17)
%
% This template was downloaded from:
% http://www.LaTeXTemplates.com
%
% Version 2.x major modifications by:
% Vel (vel@latextemplates.com)
%
% This template is based on a template by:
% Steve Gunn (http://users.ecs.soton.ac.uk/srg/softwaretools/document/templates/)
% Sunil Patel (http://www.sunilpatel.co.uk/thesis-template/)
%
% Template license:
% CC BY-NC-SA 3.0 (http://creativecommons.org/licenses/by-nc-sa/3.0/)
%
%%%%%%%%%%%%%%%%%%%%%%%%%%%%%%%%%%%%%%%%%

%----------------------------------------------------------------------------------------
%	PACKAGES AND OTHER DOCUMENT CONFIGURATIONS
%----------------------------------------------------------------------------------------

\documentclass[
11pt, % The default document font size, options: 10pt, 11pt, 12pt
%oneside, % Two side (alternating margins) for binding by default, uncomment to switch to one side
english, % ngerman for German
singlespacing, % Single line spacing, alternatives: onehalfspacing or doublespacing
%draft, % Uncomment to enable draft mode (no pictures, no links, overfull hboxes indicated)
%nolistspacing, % If the document is onehalfspacing or doublespacing, uncomment this to set spacing in lists to single
%liststotoc, % Uncomment to add the list of figures/tables/etc to the table of contents
%toctotoc, % Uncomment to add the main table of contents to the table of contents
%parskip, % Uncomment to add space between paragraphs
%nohyperref, % Uncomment to not load the hyperref package
headsepline, % Uncomment to get a line under the header
%chapterinoneline, % Uncomment to place the chapter title next to the number on one line
%consistentlayout, % Uncomment to change the layout of the declaration, abstract and acknowledgements pages to match the default layout
]{MastersDoctoralThesis} % The class file specifying the document structure

\usepackage[utf8]{inputenc} % Required for inputting international characters
\usepackage[T1]{fontenc} % Output font encoding for international characters
\usepackage{todonotes}
\usepackage{mathpazo} % Use the Palatino font by default
\usepackage{amsmath}
\usepackage{multirow}
\usepackage{multicol}
\usepackage{soul}
\usepackage{diagbox}
% \usepackage{pgfplots}
% \usepgfplotslibrary{external}
% \tikzexternalize[prefix=tikz/]

\DeclareMathOperator*{\argmax}{\arg\!\max}
\DeclareMathOperator*{\argmin}{\arg\!\min}

\usepackage[backend=bibtex,style=authoryear,natbib=true]{biblatex} % Use the bibtex backend with the authoryear citation style (which resembles APA)

\addbibresource{example.bib} % The filename of the bibliography

\usepackage[autostyle=true]{csquotes} % Required to generate language-dependent quotes in the bibliography

\newcommand{\sst}[1]{\lbl{#1}} % supersense tag label

\newcommand{\rf}[2]{\psst{#1}$\leadsto$\psst{#2}}

\definecolor{mdgreen}{rgb}{0.05,0.6,0.05}

\newcommand{\psst}[1]{\textcolor{mdgreen}{\sst{#1}}}
\newcommand{\p}[1]{\textbf{\textsf{#1}}} % preposition type

% preposition supersense tag label

%----------------------------------------------------------------------------------------
%	MARGIN SETTINGS
%----------------------------------------------------------------------------------------

\geometry{
	paper=a4paper, % Change to letterpaper for US letter
	inner=2.5cm, % Inner margin
	outer=3.8cm, % Outer margin
	bindingoffset=.5cm, % Binding offset
	top=1.5cm, % Top margin
	bottom=1.5cm, % Bottom margin
	%showframe, % Uncomment to show how the type block is set on the page
}

%----------------------------------------------------------------------------------------
%	THESIS INFORMATION
%----------------------------------------------------------------------------------------

\thesistitle{Preposition Supersenses: Disambiguation, Application, and Analysis} % Your thesis title, this is used in the title and abstract, print it elsewhere with \ttitle
\supervisor{Dr. Omri \textsc{Abend}} % Your supervisor's name, this is used in the title page, print it elsewhere with \supname
\examiner{} % Your examiner's name, this is not currently used anywhere in the template, print it elsewhere with \examname
\degree{Msc. of Computer Science} % Your degree name, this is used in the title page and abstract, print it elsewhere with \degreename
\author{Aviram \textsc{Stern}} % Your name, this is used in the title page and abstract, print it elsewhere with \authorname
\addresses{} % Your address, this is not currently used anywhere in the template, print it elsewhere with \addressname

\subject{Preposition Supersenses} % Your subject area, this is not currently used anywhere in the template, print it elsewhere with \subjectname
\keywords{} % Keywords for your thesis, this is not currently used anywhere in the template, print it elsewhere with \keywordnames
\university{\href{http://www.huji.ac.il}{The Hewbrew University of Jerusalem}} % Your university's name and URL, this is used in the title page and abstract, print it elsewhere with \univname
\department{\href{https://www.cs.huji.ac.il/}{Benin School of Computer Science and Engineering}} % Your department's name and URL, this is used in the title page and abstract, print it elsewhere with \deptname
\group{{Computational Linguistics and Natural Language Processing lab}} % Your research group's name and URL, this is used in the title page, print it elsewhere with \groupname
\faculty{\href{https://www.cs.huji.ac.il/}{Benin School of Computer Science and Engineering}} % Your faculty's name and URL, this is used in the title page and abstract, print it elsewhere with \facname

\AtBeginDocument{
\hypersetup{pdftitle=\ttitle} % Set the PDF's title to your title
\hypersetup{pdfauthor=\authorname} % Set the PDF's author to your name
\hypersetup{pdfkeywords=\keywordnames} % Set the PDF's keywords to your keywords
}
\usepackage{svg}
\usepackage[edges]{forest}
\usepackage{xytree}
\usepackage{amssymb}
\usepackage{gb4e} \noautomath

\begin{document}

\frontmatter % Use roman page numbering style (i, ii, iii, iv...) for the pre-content pages

\pagestyle{plain} % Default to the plain heading style until the thesis style is called for the body content

%----------------------------------------------------------------------------------------
%	TITLE PAGE
%----------------------------------------------------------------------------------------

\begin{titlepage}
\begin{center}

\vspace*{.06\textheight}
{\scshape\LARGE \univname\par}\vspace{1.5cm} % University name
\textsc{\Large Masters Thesis}\\[0.5cm] % Thesis type

\HRule \\[0.4cm] % Horizontal line
{\huge \bfseries \ttitle\par}\vspace{0.4cm} % Thesis title
\HRule \\[1.5cm] % Horizontal line
 
\begin{minipage}[t]{0.4\textwidth}
\begin{flushleft} \large
\emph{Author:}\\
{\authorname} % Author name - remove the \href bracket to remove the link
\end{flushleft}
\end{minipage}
\begin{minipage}[t]{0.4\textwidth}
\begin{flushright} \large
\emph{Supervisor:} \\
\href{http://www.cs.huji.ac.il/~oabend/}{\supname} % Supervisor name - remove the \href bracket to remove the link  
\end{flushright}
\end{minipage}\\[3cm]
 
\vfill

\large \textit{A thesis submitted in fulfillment of the requirements\\ for the degree of \degreename}\\[0.3cm] % University requirement text
\textit{in the}\\[0.4cm]
\groupname\\\deptname\\[2cm] % Research group name and department name
 
\vfill

{\large \today}\\[4cm] % Date
%\includegraphics{Logo} % University/department logo - uncomment to place it
 
\vfill
\end{center}
\end{titlepage}

%----------------------------------------------------------------------------------------
%	DECLARATION PAGE
%----------------------------------------------------------------------------------------

\begin{declaration}
\addchaptertocentry{\authorshipname} % Add the declaration to the table of contents
\noindent I, \authorname, declare that this thesis titled, \enquote{\ttitle} and the work presented in it are my own. I confirm that:

\begin{itemize} 
\item This work was done wholly or mainly while in candidature for a research degree at this University.
\item Where any part of this thesis has previously been submitted for a degree or any other qualification at this University or any other institution, this has been clearly stated.
\item Where I have consulted the published work of others, this is always clearly attributed.
\item Where I have quoted from the work of others, the source is always given. With the exception of such quotations, this thesis is entirely my own work.
\item I have acknowledged all main sources of help.
\item Where the thesis is based on work done by myself jointly with others, I have made clear exactly what was done by others and what I have contributed myself.\\
\end{itemize}
 
\noindent Signed:\\
\rule[0.5em]{25em}{0.5pt} % This prints a line for the signature
 
\noindent Date:\\
\rule[0.5em]{25em}{0.5pt} % This prints a line to write the date
\end{declaration}

\cleardoublepage

% %----------------------------------------------------------------------------------------
% %	QUOTATION PAGE
% %----------------------------------------------------------------------------------------

% \vspace*{0.2\textheight}

% \noindent\enquote{\itshape Thanks to my solid academic training, today I can write hundreds of words on virtually any topic without possessing a shred of information, which is how I got a good job in journalism.}\bigbreak

% \hfill Dave Barry

%----------------------------------------------------------------------------------------
%	ABSTRACT PAGE
%----------------------------------------------------------------------------------------

\begin{abstract}
\addchaptertocentry{\abstractname} % Add the abstract to the table of contents
% The Thesis Abstract is written here (and usually kept to just this page). The page is kept centered vertically so can expand into the blank space above the title too\ldots
Function words such as adpositions and possessives are crucial for conveying the semantic relations between different parts of a sentence. This work focuses on a recently introduced schema attempting to characterize these relations - \textbf{Semantic Network of Adposition and Case Supersense} (SNACS) by (\cite{snacs}). The schema consists of 50 labels - \textit{Preposition Supersenses} (PSS) - each representing some class of semantic relations attributed to a preposition in a sentence. 

Chapter 1 deals with supervised methods for automatic disambiguation of preposition supersenses. We present a neural PSS model, trained and evaluated on the STREUSLE corpus of PSS annotated English web reviews (\cite{snacs}). We experiment with several settings of model architecture, auxiliary information acquisition and target identification. Specifically, we find that a model using contextualized word embeddings significantly outperforms a model using non-contextualized word embeddings with auxiliary syntactic information, as previously demonstrated by (\cite{nelson}). 

Chapter 2 examines how beneficial the information provided by PSS is to the Prepositional-Phrase (PP) Attachment problem - the task of resolving the ambiguity of which phrase is referred to by a prepositional phrase, as demonstrated by the following sentence: “I [saw]VP [the man]NP [with the telescope]PP”. The prepositional phrase - \textit{“with the telescope”} may acts as a modifier to the verb \textit{saw} (the speaker used the telescope), or to the noun phrase - \textit{"the man"} (the man had the telescope). In such scenarios, each interpretation may yield a different preposition supersense for the corresponding preposition, thus providing evidence to believe that a PP-Attachment disambiguator might benefit from using PSS information. We experiment with an existing PP-Attachment model, by incorporating PSS features onto it, expecting to see an improvement in disambiguation accuracy. Empirical results show a small, possibly insignificant improvement in accuracy of the PSS-enhanced model. Another observation obtained by experimenting with several datasets indicates that PSS information behaves as a normalizing factor between different domains of content.

Chapter 3 examines the transferability of the knowledge acquired by a PSS
disambiguator between different languages and domains, by evaluating
the PSS model presented in chapter \ref{Chapter2} - trained solely using the STREUSLE corpus of English web reviews - on the recently introduced Mandarine Chinese PSS corpus
by (\cite{chinesecorpus}), with the use of parallel text and word alignment models. We find that domain and language have a critical impact on model peformence, as the accuracies on the Chinese corpus drop to a near-baseline level.

% \include{Chapters\Abstract.tex}
\end{abstract}

%----------------------------------------------------------------------------------------
%	ACKNOWLEDGEMENTS
%----------------------------------------------------------------------------------------

\begin{acknowledgements}
\addchaptertocentry{\acknowledgementname} % Add the acknowledgements to the table of contents
% The acknowledgments and the people to thank go here, don't forget to include your project advisor\ldots
\todo{Write acknowledgments}
\end{acknowledgements}

%----------------------------------------------------------------------------------------
%	LIST OF CONTENTS/FIGURES/TABLES PAGES
%----------------------------------------------------------------------------------------

\tableofcontents % Prints the main table of contents

\listoffigures % Prints the list of figures

\listoftables % Prints the list of tables

%----------------------------------------------------------------------------------------
%	ABBREVIATIONS
%----------------------------------------------------------------------------------------

\begin{abbreviations}{ll} % Include a list of abbreviations (a table of two columns)

% \textbf{LAH} & \textbf{L}ist \textbf{A}bbreviations \textbf{H}ere\\
\textbf{PSS} & \textbf{P}reposition \textbf{S}uper\textbf{s}ens\\
\textbf{PP} & \textbf{P}repositional \textbf{P}hrase\\
\textbf{NP} & \textbf{N}oun \textbf{P}hrase\\
\textbf{VP} & \textbf{V}erb \textbf{P}hrase\\
\textbf{PPA} & \textbf{P}repositional \textbf{P}hrase \textbf{A}ttachment\\

\end{abbreviations}

% %----------------------------------------------------------------------------------------
% %	PHYSICAL CONSTANTS/OTHER DEFINITIONS
% %----------------------------------------------------------------------------------------

% \begin{constants}{lr@{${}={}$}l} % The list of physical constants is a three column table

% % The \SI{}{} command is provided by the siunitx package, see its documentation for instructions on how to use it

% Speed of Light & $c_{0}$ & \SI{2.99792458e8}{\meter\per\second} (exact)\\
% %Constant Name & $Symbol$ & $Constant Value$ with units\\

% \end{constants}

% %----------------------------------------------------------------------------------------
% %	SYMBOLS
% %----------------------------------------------------------------------------------------

% \begin{symbols}{lll} % Include a list of Symbols (a three column table)

% $a$ & distance & \si{\meter} \\
% $P$ & power & \si{\watt} (\si{\joule\per\second}) \\
% %Symbol & Name & Unit \\

% \addlinespace % Gap to separate the Roman symbols from the Greek

% $\omega$ & angular frequency & \si{\radian} \\

% \end{symbols}

%----------------------------------------------------------------------------------------
%	DEDICATION
%----------------------------------------------------------------------------------------

\dedicatory{For/Dedicated to/To my\ldots} 
\todo{Write dedication}

%----------------------------------------------------------------------------------------
%	THESIS CONTENT - CHAPTERS
%----------------------------------------------------------------------------------------

\mainmatter % Begin numeric (1,2,3...) page numbering

\pagestyle{thesis} % Return the page headers back to the "thesis" style

% Include the chapters of the thesis as separate files from the Chapters folder
% Uncomment the lines as you write the chapters

% Chapter Template

\chapter{Background} % Main chapter title

\label{Background} % Change X to a consecutive number; for referencing this chapter elsewhere, use \ref{ChapterX}


%-----------------------------------
%	SECTION 1
%-----------------------------------
\section{Introduction to Preposition Supersenses}
\subsubsection{*Adpositions signal semantic relations}
\subsubsection{*These relations are ambiguous}
\subsubsection{*Disambiguations could be beneficial to downstream tasks}
\subsection{The SNACS Scheme}
\subsubsection{SNACS Label Hierarchy}
\subsubsection{Construal Analysis}
\subsection{The STREUSLE corpus}
\subsection{Previous Approaches}
\subsection{PSS in This Work}
\pagebreak

%-----------------------------------
%	SECTION 2
%-----------------------------------
\section{Neural Networks for NLP}
\subsubsection{*Neural nets overview}
\subsubsection{*Activation functions}
\subsubsection{*Loss functions}
\subsubsection{*Training}
\subsubsection{*Hyperparameters vs. Model parameters, tuning}
\subsubsection{*Regularization and drop-out}
\subsection{Feed Forward Networks}
\subsection{Recurrent Neural Networks}
\subsubsection{Bidirectional RNN}
\subsubsection{Stacked RNNs}
\subsubsection{Long-Short-Term-Memory (LSTM)}
\subsection{Lookup Tables}
\subsection{The Stacked BiLSTM-MLP-Softmax Architecture}
\subsection{Neural Networks in this work}
\pagebreak

%-----------------------------------
%	SECTION 3
%-----------------------------------
\section{Word Embeddings Overview}
\subsection{Main Approches}
\subsubsection{The Skip-Gram Model (Word2Vec)}
\subsubsection{Global Vectors for Word Representation (GLoVe)}
\subsubsection{Syntactic Embeddings}
\subsubsection{FastText?}
\subsection{Contextualized Embeddings}
\subsubsection{ELMo}
\subsection{Multi-Lingual Word Embeddings}
\subsubsection{MUSE}
\subsection{Learning via Word Embeddings}
\subsubsection{*Basic usage with a FF-Network}
\subsubsection{*Updating embeddings through training}
\subsubsection{*Contextutalized representations with RNNs}
\subsection{Handling Out-Of-Vocabulary words}
\subsubsection{*By stemming}
\subsubsection{*By learning an internal representation}
\subsubsection{*Character-based embeddings}
\subsection{Use of Word Embeddings in this work}
\pagebreak


%-----------------------------------
%	SECTION 4
%-----------------------------------
\section{Natural Language Parsing and the Prepositional Phrase Attachment Problem}
\subsubsection{*Outline}
\subsection{Sentence Parsing}
\subsubsection{Dependency vs. Constituency Based Parsing}
\subsubsection{Syntactic vs. Semantic Parsers}
\subsection{The PP-Attachment problem}
\subsubsection{*Overview, examples}
\subsection{PP Attachment Disambiguation Approaches}
\subsubsection{As a sub-task of sentence parsing}
\subsubsection{As an isolated task}
\subsection{PP Attachment and PSS}
\subsubsection{*Motivation of using PSS for PPA}
\pagebreak

%-----------------------------------
%	SECTION 5
%-----------------------------------
\section{PPA work by Belinkov et. al 2014}
\subsubsection{*Overview, relevance to this work}
\subsubsection{*Intro to Belinkov work}
\subsubsection{* -- Isolated PPA models}
\subsubsection{* -- Incorporated to an existing parser}
\subsubsection{* -- Model architecture in brief}
\subsection{Datasets}
\subsubsection{*Structure, size and other notable stats}
\subsubsection{*Extraction process}
\subsection{Model}
\subsubsection{*Model framework}
\subsubsection{*Score function variations}
\subsubsection{The HPCD Variation}
\subsection{Word Vector Representation}
\subsubsection{*Word Embedding Variations}
\subsubsection{*Vector Enrichment (WordNet, VerbNet, etc)}
\subsubsection{*Full Vector Outline}
\subsection{Training}
\subsection{Results}
\subsection{Conclusion}
\subsubsection{*Brief summary and chapter 3 outline}
\pagebreak

%-----------------------------------
%	SECTION 6
%-----------------------------------
\section{Universal Dependencies}
\pagebreak

%-----------------------------------
%	SECTION 7
%-----------------------------------
\section{Named Entity Recognition}

% Chapter Template

\chapter{Preposition Supersenses Disambiguation} % Main chapter title

\label{Chapter2} % Change X to a consecutive number; for referencing this chapter elsewhere, use \ref{ChapterX}

\section{Overview}
\textbf{Semantic Network of Adposition and Case Supersense} (SNACS), by (\cite{snacs}), is a labeling schema assigning semantic attributes to adpositions. Each adposition is assigned with a pair of \emph{role} and \emph{function} supersenses, from a bank of 50 labels called \emph{Preposition Supersenses} or \emph{PSS} \footnote{See section \ref{sec:snacsscheme} for a detailed description of the SNACS schema.}. These annotations provide an additional layer of semantic information regarding a sentence, which could be beneficial for downstream NLP applications. In this chapter, we present a neural PSS disambiguation model, trained and evaluated on the STREUSLE corpus of gold-standard PSS annotations. 

This chapter begins with a more formal description of the task as a pipline of two subtasks - (i) target identification (section \ref{sec:probform_ident}) - identifying the adpositions of a sentence requiring a PSS annotation, and (ii) disambiguation (section \ref{sec:probform_disam}) - deciding on the correct PSS annotation of an identified adposition. It follows with an in-depth description of the STREUSLE corpus used for training and evaluation (section \ref{sec:streusle}), and the heuristic process used for target identification (section \ref{sec:ident}). The following sections present the architecture of the neural PSS disambiguation model and data preprocessing steps (section \ref{sec:pssdisambg}), the training process (section \ref{sec:psstraining}), and the different model and evaluation settings used throughout the experiment (section \ref{sec:psssettings}).  Finally, the evaluation of the model compared to a set of baselines is presented and analysed (sections \ref{sec:pssbaselines}, \ref{sec:pssresults}, and \ref{sec:pssanalysis}).

\section{Problem Formulation} \label{sec:probform}
\subsection{Target Identification}\label{sec:probform_ident}
The target identification task deals with identifying the PSS-bearing units in a given sentence. A PSS bearing unit is not limited to just prepositions (despite the confusing terminology) - a PSS can also be applied to postpositions (\emph{"ago", "aside"}), cases (\emph{"Harry\textbf{'s} wand", "\textbf{his} wand"}) and many other phrases that impose a preposition-like behaviour (\emph{"in front of", "out of"})\footnote{A more thorough categorization of PSS bearing units is provided in section \ref{sec:snacsscheme}}.  The set of PSS bearing units is therefore comprised of non-overlapping, consecutive token sequences, as defined by the SNACS specification and annotation guide. 

Most PSS-bearing units are plain adpositions, and thus can be easily identified using POS labels or even a predefined adposition vocabulary. In many cases, however, lexicality alone is not enough to determine a preposition. The word \emph{to}, for instance, is considered a PSS-bearing preposition in \emph{"They worked \textbf{to} fix the issue."} (PURPOSE), but not considered one in \emph{"I need \textbf{to} leave."} (infinitive-clause), despite the similar structure. POS information is also not enough for a comprehensive identification - while 61.2\% of the units annotated in the STREUSLE corpus are adpositions according to gold POS annotation, 20.2\% are possessives, and 18.6\% belong
to other POS classes. The task of identifying the correct PSS units is therefore not trivial as one might expect. This work makes use of a PSS identification heuristic by (\cite{snacs}), described in section \ref{sec:ident}.

\subsection{Supersense Disambiguation}\label{sec:probform_disam}

The identified units are to be assigned with a \emph{Scene Role} and \emph{Function} PSS, both from the same bank of labels defined in the SNACS specification (section \ref{sec:construal}). Role and Function labels tend to collide, as demonstrated by the STREUSEL corpus, where 67.8\% of instances are annotated with the same label for both. The neural PSS disambiguation model presented in section \ref{sec:pssdisambg} jointly predicts both labels.

\section{The STREUSLE corpus}\label{sec:streusle}

\begin{figure}
\begin{exe}\raggedright
\ex This cottage is a charming homely, friendly, place \p{to}/\rf{Characteristic}{Purpose} stay.
\ex I came in \p{to}/\psst{Purpose} get a nice gift \p{for}/\rf{Recipient}{Direction} \p{my}/\rf{SocialRel}{Gestalt} wife.
\ex He worked \p{on}/\psst{Theme} it right \p{on}/\psst{Locus} the back \p{of}/\psst{Whole} \p{my}/\psst{Possessor} car.
\end{exe}
\caption{Annotated sentences from the STREUSLE corpus. PSS labels are annotated as \rf{Role}{Function}, or a single lable when they collide}
\label{fig:streuslesamples}
\end{figure}


The STREUSLE corpus of PSS annotated sentences (\cite{snacs}) was constructed by applying the SNACS scheme to a collection of online consumer reviews taken from the English Web Treebank (\cite{ewtb}). Standard train/dev/test splits are used, consisting of 2,723/554/535 sentences and 4,522/453/480 annotated targets respectively. 305 (5.5\%) of units  are multiword expressions, and for a total of 3,702 (67.8\%) units the Role and Function labels collide. The corpus contains 249 different types of PSS-bearing units, from which 149 are mutiword expressions. The most frequent PSS label in the corpus is LOCUS - annotated for Role (11.6\%), Function (14.2\%), and both (10.4\%).

Sentences in the English Web Treebank corpus are equipped with manually -annotated syntactic information, in the form of tokenization, lemmatization, and a full Universal Dependencies parse tree. The STREUSLE corpus is bundled with this information, which from here on out is referred to as \emph{gold syntax}. Similarly, the set of manually-identified PSS bearing units are refered to as \emph{gold identification}. Example sentences from the corpus, along with their PSS annotations, are available in figure \ref{fig:streuslesamples}.

\section{Target Identification Heuristic} \label{sec:ident}

A fully automated PSS annotation system requires the ability to automatically identify PSS-bearing units, without relying on human annotated targets such as in the STREUSLE corpus. This work makes use of a PSS target identification heuristic by (\cite{snacs}), jointly developed along with the neural PSS disambiguator presented in this work in a collaborative manner, to facilitate a fully automated PSS annotator. 

As a preprocקssing step, the algorithm makes use of the training portion of the STREUSLE corpus to construct a set of lexical lists:
\begin{enumerate}
    \item A blacklist of non PSS-bearing multiword expressions containing preposition tokens (e.g., \emph{take care of}).
    \item A whitelist of PSS-bearing multiword expressions.
    \item A whitelist of PSS-bearing single word expressions.
    \item A whitelist and a blacklist of governors for PSS-bearing and non PSS-bearing instances of \emph{to} respectively, 
\end{enumerate}

For (4), the algorithm makes use of the UD parse-tree. List (1), (2) and (3) are enriched with a set of manually identified lexical terms. The identification process is then performed as follows: sentences are first scanned for multiword expressions using lists (1) and (2). From the remaining segments, single-word candidates are identified by matching a high-recall set of POS tags,
then filtered through 5 different heuristics for adpositions, possessives, subordinating conjunctions,
adverbs, and infinitivals. Most heuristics use a combination of rules involving syntax information and the use of list (3) for filtering. The infinitival heuristic also makes use of the learned whitelist and blacklist described in (4).

\begin{table}[t]\centering\small
\begin{tabular}{lccc}
Syntax & P & R & F \\
\toprule
gold & 88.8 & 89.6 & 89.2 \\
auto & 86.0 & 85.8 & 85.9 \\
\end{tabular}
\caption{Target identification results on the STREUSLE test set, when using gold vs. automatically parsed syntax information.}
\label{tab:targetid}
\end{table}

Table \ref{table:identresults}
reports the precision, recall, and F-score (P/R/F) of
the target identification heuristic on the test portion of the STREUSLE corpus. This evaluation compares the use of gold vs. automatically parsed syntax information, showing a slight decrease of quality when using the latter. 47 of 54 false positives (87\%) can be ascribed to tokens that are part of a (non-adpositional or larger adpositional) multiword expression. 9 of the 50 false negatives (18\%) are rare multiword expressions not
occurring in the training data, and there are 7 partially identified ones, which are counted as both
false positives and false negatives.

\section{PSS Disambiguation Model} \label{sec:pssdisambg}

\begin{figure}
    \centering
    \includegraphics{Figures/Electron.pdf}
    \caption{The full neural PSS dismabiguation model. Tokens are fed into stacked BiLSTM layers. For each PSS-bearing unit, a feature vector is constructed and fed into two MLP-Softmax layers, predicting label probabilities for Role and Function}
    \label{fig:pssmodel}
\end{figure}

The neural PSS disambiguation model presented in this work is a multilayer perceptron (MLP) stacked on top of a multilayered BiLSTM. Such architectures have proven to be successful for a variety of NLP tasks (\cite{kiperwasser16simple, elmo}). The model architecture is largely comparable to that of (\cite{gonen16semi}), who experimented with a coarsened version of STREUSLE. The main difference is their use of unlabeled multilingual datasets to improve prediction by exploiting the differences in preposition ambiguities across languages. The model makes use of syntactic information in the form of tokenization, lemmatization, NER annotation, and a full Universal Dependencies parse tree, all which are assume to be given as inputs in addition to the actual sentence and target words.  The full model architecture is depicted in figure \ref{fig:pssmodel}. 

Each sentence word $w_i$ is first converted into a vector using pretrained word embeddings, and a vocabulary-sized lookup table for an internal representation of each word. We experiment with contextualized (ELMo) and non-contextualized (SkipGram) pretrained embeddings\footnote{See section \ref{section:wordembeddings} for more details regarding each embedding method.}. For ELMo, the pretrained model from the AllenNLP python module (\cite{allennlp}) is used. For skip gram, the pretrained FastText vectors are used (\cite{fasttext}). The full word embedding $e_i$ is therefore a concatenation of three vectors: the pretrained embeddings of the word $e^{word}_i$ and lemma $e^{lemma}_i$, and the internal word representation $e^{internal}_i$. 

$$ e_i = [e^{word}_i;e^{lemma}_i;e^{internal}_i] $$

Hyperparameters control which pretrained embeddings to use (ELMo/SkipGram), and in the case of fastText, whether the pretrained vectors remain fixed or get updated during training. A single ELMo layer in the pretrained model consists of 1000 elements, whereas for fastText the vectors are of size 300. The dimension of the internal representation (if used) is also determined by a hyperparameter. The size of the full vector is therefore largely dependent of the choice of hyperparameters. A $NONE$ label is used to represent words lacking a pretrained representation, both in training and at test time.

The resulting word embeddings $e_1, .. ,e_n$ are then fed into a multilayered BiLSTM, yielding contextualized word representations $c_1, .. ,c_n$. The number of LSTM layers, the dimension of the LSTM output, and whether the bidirectional variant is used, are all controlled by hyperparameters.

$$ c_1, ... , c_n = Multilayerd\_BiLSTM(e_1, ... ,e_n)$$

A vector $u_{start:end}$ is then constructed for an identified PSS-bearing target span $(w_{start}, .., w_{end})$, using the corresponding contextualized vectors and the auxiliary syntactic information (collectively denoted here as \emph{syntax}). The exact structure of this vector, which is refereed to as the Preposition Feature Vector, is explained in section \ref{sec:enrichedvec}. 

$$u_{start:end} = Build\_Prep\_Feature\_Vec(c_1, ... ,c_n, start, end, syntax)$$

Each feature vector $u$ is then fed into two separate MLPs, followed by separate Softmax layers that
yield the predicted probabilities for the Role and
Function labels:

$$ p^{Role}_{start:end} = Softmax(MLP_{Role}(u_{start:end})) $$
$$ p^{Function}_{start:end} = Softmax(MLP_{Function}(u_{start:end})) $$

$p^{Role}_{start:end}$ and $p^{Function}_{start:end}$ are then the corresponding predicted probability distributions for Role and Function labels respectively. The number of MLP layers and dimensions, as well as the activation functions used by the MLP layers, are all considered model hyperparameters. 

\subsection{Constructing the Preposition Feature Vector} \label{sec:enrichedvec}
The preposition feature vector is obtained by constructing and concatenating feature vectors for the preposition, and its object and governor headwords. The object and governor of a preposition are determined in a heuristic process (described in section \ref{sec:govobj}), which also emits the preposition's lexical category - a syntactic label indicating whether the preposition is
predicative/stranded/possesive/subordinating/none of these. 

A single feature vector for a word $w_i$ is obtained by concatenating its contextualized representation $c_i$, with the embeddings of its POS, UD dependency and NER labels. The preposition vector is also concatenated with the embedding of its lexical category, and a binary indicator of whether either of the two tokens following the preposition is capitalized. In the case of multiword prepositions, the first token is used. All labels are embedded using randomly initialized lookup tables, with dimensions determined by hyperparameters.

\subsection{Extracting the Governor and Object of a Preposition} \label{sec:govobj}

The governor and object of a preposition may provide a great deal of information regarding the semantics of the preposition. The heuristic process used for identifying the governor and object of a preposition, by (\cite{snacs}) was developed as part of a collaborative effort on the task of PSS disambiguation. In addition to identification, it also sorts the prepositions into 5 lexical categories: stranded (\emph{"The book I told you \textbf{about}"}), predicative (\emph{"The dog is \textbf{under} the table"}), possessive (\emph{"The dog\textbf{'s} food"}), subordinating (\emph{"I will arrive \textbf{before} sunset"}), or none of these. The heuristic applies a set of rules using the UD parse tree and POS labels: first, the UD dependency label of the preposition is examined. For prepositions with \emph{case} or \emph{mark}, the parent and grandparent of the preposition in the UD tree are usually the object and governor, respectively. Prepositions with different labels are typically not attached to an object, with their parent being the governor.  The \emph{nmod:poss} UD label is an indication of a possessive preposition. The heuristic process then attempts to handle exceptions to the above, by explicitly matching against patterns corresponding to stranded prepositions and copular constructions.  
\section{Training the Model} \label{sec:psstraining}

The model was trained using the cross-entropy loss function. More specifically, the loss of a single identified PSS target is defined as the negated log-sum of the probabilities assigned to the correct labels: 
$$ Loss(x, y) = - log(p_r) - log(p_f) $$

Where $p_r$ and $p_f$ are the probabilities assigned by the model to the correct Role and Function labels ($y$) for the given sentence and identified PSS unit ($x$).

Model parameters are initialized randomly using the Xavier initialization (\cite{glorot10xavier}). Inverted dropout was applied for the LSTM and MLP layers. Training was performed using SGD with Backpropagation for 130 epochs, and a minibatch of size 20. To avoid overfitting, the model parameters corresponding to the best performing epoch on the development set are selected. The dropout probabilities, SGD learning rate and learning rate decay factor are all hyperparameters. The model was implemented with the DyNet library (\cite{dynet}).

\subsubsection{Hyperparameter Tuning}

The development section of the STREUSLE corpus was used to tune hyperparameters using a random search process. Hyperparameters are iteratively sampled and used to train a model in the process described above. The best performing set of hyperparameters is selected. Hyperparameters were used in this manner for tuning the structure of the network (number of layers, dimensions, activations), tuning the optimizer (learning rate and learning rate decay), and feature selection (e.g. which syntactic information and word embeddings to use). The complete list of hyperparameters is available in table (\ref{tab:hyperparams}).

\section{Experiment Settings} \label{sec:psssettings}

The aformetioned PSS dismabiguation model heavily relies on two preprocessing steps: syntactic information extraction (tokenization, lemmatization, UD parsing, NER labeling, and so on), followed by PSS target identification (section \ref{sec:probform_ident}). The required syntactic information - with the exception of NER labels - is entirely available through the gold standard annotations of the STREUSLE corpus, referred to here as \textit{gold-syntax}. To better handle real-life scenarios where such syntactic information is not available, the model is also evaluated using an automatically parsed syntactic information instead, by applying the CoreNLP parser (\cite{manning14stanford}) to the entire corpus. The settings in which such syntactic information is used are referred to as \textit{auto-syntax}. Each setting is using the corresponding syntactic information source throughout the entire process of target identification, training and testing. 

We take a similar approach with the target identification step: the model is evaluated using the gold-standard identifications in the STREUSLE corpus (\textit{gold-id} settings), and also using the automatic identification heuristic from section \ref{sec:ident} (\textit{auto-id} settings).

The model is evaluated for each of the four combinations of \textit{gold-id/auto-id} and \textit{gold-syntax/auto-syntax}. The training and tuning processes described above are applied to each setting individually.

\section{Baselines} \label{sec:pssbaselines}

In addition to the neural model, we also evaluate the following two PSS disambiguation baselines:

\begin{itemize}
    \item \textbf{Most Frequent Label (MF)} - always predicting the most frequent PSS label seen in the training data.
    \item \textbf{Most Frequent Label per Preposition (MF-Prep)} - predicting the most frequent PSS label for the specific (lemmatized) preposition, according to the training data.
\end{itemize}

Both baselines are evaluated and compared with the full model for each of the four experiment settings. 

\section{Results} \label{sec:pssresults}

\begin{table*}[]\label{tab:pssresults}
    \newcommand{\score}[2]{#1{\tiny ($\pm$#2)}}
    \newcommand{\acc}[1]{\multicolumn{3}{c|}{\textcolor{gray}{\rule[2pt]{0.5in}{0.5pt}} #1 \textcolor{gray}{\rule[2pt]{0.5in}{0.5pt}}}}
    \newcommand{\accl}[1]{\multicolumn{3}{c}{\textcolor{gray}{\rule[2pt]{0.5in}{0.5pt}} #1 \textcolor{gray}{\rule[2pt]{0.5in}{0.5pt}}}}
    \newcommand{\accs}[1]{\multicolumn{3}{c|}{\textcolor{gray}{\rule[2pt]{0.3in}{0.5pt}} #1 \textcolor{gray}{\rule[2pt]{0.3in}{0.5pt}}}}
    \newcommand{\accsl}[1]{\multicolumn{3}{c}{\textcolor{gray}{\rule[2pt]{0.3in}{0.5pt}} #1 \textcolor{gray}{\rule[2pt]{0.3in}{0.5pt}}}}

	\centering\footnotesize
    \setlength{\tabcolsep}{1pt} % Default value: 6pt
    \renewcommand{\arraystretch}{1} % Default value: 1
	\begin{tabular}{@{}ccc<{\hspace{5pt}}ccc|ccc|ccc@{}}
% 		\toprule
                          & &          & \multicolumn{3}{c}{\textit{Role}} & \multicolumn{3}{c}{\textit{Func.}} & \multicolumn{3}{c}{\textit{Full}}              \\
                          & \textbf{Syntax} & \textbf{ID}       & P                                 & R                                  & F    & P    & R    & F    & P    & R    & F    \\
		\midrule
        \multirow{2}{*}{MF} & \multirow{2}{*}{N/A} & gold & \acc{ 15.8 } & \acc{ 19.8 } & \accl{ 14.8 } \\
          &   & auto & 14.3 & 14.4 & 14.3 & 16.5 & 16.7 & 16.6 & 13.2 & 13.3 & 13.3 \\
        \midrule
        \multirow{2}{*}{MF-Prep} & \multirow{2}{*}{N/A} & gold & \acc{ 40.6 } & \acc{ 53.3 } & \accl{ 37.9 } \\
          &   & auto & 37.0 & 37.3 & 37.1 & 49.8 & 50.2 & 50.0 & 34.3 & 34.6 & 34.4 \\        \midrule
        \multirow{4}{*}{Neural}        & \multirow{2}{*}{gold} & gold & \accs{\score{ 79.5 }{ 0.6 } }& \accs{\score{ 86.8 }{ 0.9 } }& \accsl{\score{ 75.7 }{ 0.6 } }\\
         &   & auto & \score{ 68.6 }{ 1.0 } & \score{ 69.2 }{ 1.0 } & \score{ 68.9 }{ 1.0 } & \score{ 75.5 }{ 0.6 } & \score{ 76.1 }{ 0.6 } & \score{ 75.8 }{ 0.6 } & \score{ 65.2 }{ 1.0 } & \score{ 65.7 }{ 1.0 } & \score{ 65.4 }{ 1.0 } \\
         \cmidrule{2-12}                & \multirow{2}{*}{auto} & gold & \accs{\score{ 77.5 }{ 1.2 } } & \accs{\score{ 85.9 }{ 0.9 } }& \accsl{\score{ 74.0 }{ 1.4 } }\\
         &   & auto & \score{ 66.4 }{ 0.8 } & \score{ 66.2 }{ 0.9 } & \score{ 66.3 }{ 0.9 } & \score{ 74.9 }{ 0.6 } & \score{ 74.8 }{ 0.5 } & \score{ 74.8 }{ 0.6 } & \score{ 63.6 }{ 0.6 } & \score{ 63.5 }{ 0.6 } & \score{ 63.6 }{ 0.6 } \\
         \midrule
% 		\bottomrule
	\end{tabular}
	\caption{\label{tab:overall} Overall performance of SNACS disambiguation systems on the test set. Results are reported for the role supersense ({\it Role}), the function supersense ({\it Func.}), and their conjunction ({\it Full}). The average and standard deviation for 5 different training instances are reported. All figures are percentages.}
\end{table*}


The precision, recall, and $F_1$ score ($P/R/F$) for the best performing models and baselines on the test set, in each setting, are listed in table \ref{tab:pssresults}. The corresponding model hyperparameters in each setting are listed in table \ref{tab:hyperparams}. For the \textit{gold-id} settings, where $P = R = F$, a single accuracy score is reported. Due to the relatively small size of samples in the test set (480 prepositions), the scores may vary slightly but noticeably between one instance of a trained model to another (trained using different random seeds). For that reason, the average of 5 different training instances is reported, along with the standard deviation.

The neural systems outperform the most-frequent baselines by a large margin for both Role and Function label prediction across all settings. Another observation is that Function prediction is consistently more accurate then Role prediction, by a margin of roughly 10 points. This is probably due to the reduced ambiguity of functions compared to roles (as can be observed in the most-frequent baseline scores). The use of automatic syntax consistently decreases accuracy scores across all systems (a $\sim$3 point gap). In fact, the best performing model in the auto-syntax/auto-id setting does not use syntax at all (see table \ref{tab:hyperparams}). The target identification accuracy (table \ref{tab:targetid}) imposes a performance bottleneck for the auto-id settings, as can be observed by a consistent $\sim$10 point drop. 

\section{Analysis} \label{sec:pssanalysis}

\subsection{Implication of Word Embeddings and Syntax}

Contextualized word embeddings, and ELMo in particular, have consistently been shown to outperform non-contextualized embedding methods such as SkipGram and GloVe in many kinds of prediction tasks (\cite{elmo, nelson}). (\cite{nelson}) in particular demonstrate that for PSS disambiguation, a simple linear model over the ELMo representations of prepositions significantly outperforms complex and feature-rich models that make use of non-contextualized embeddings and auxiliary syntactic information. In fact, one of their baselines is a previous version (\cite{snacs}) of the PSS model described in this chapter, which used pretrained, non-contextualized  Word2Vec  embeddings\footnote{The Word2Vec model was pretrained on the Google News corpus.}. The results presented in this chapter are consistent with these findings, as can be seen by the fact that ELMo embeddings yield the best scores across all settings.

As it becomes apparent that ELMo embeddings are able to capture syntactic and semantic information, it is unclear how useful is auxiliary syntactic information when ELMo embeddings are also provided. To assess this, we evaluate 4 variations of the model, corresponding to the choice of word embeddings (ELMo/fastText) and the presence of auxiliary syntax. The rest of the hyperparameters (activations, number of layers, layer dimensions, etc.) are tuned separably for each variation. The accuracy scores for each variation in the gold-id/gold-syntax setting are reported in table \ref{tab:embdsyn}.

As expected, using the non-contextualized fastText embeddings alone yields similar results to the MF-Prep baseline, as the embeddings represent type level prepositions. Adding syntactic signals on top of fastText embeddings introduces the model with crucial contextualized information, and as a result yielding a $\sim$30 point gain in accuracy. Consistent with the findings of (\cite{nelson}), a system relying solely on contextualized ELMo embeddings (without syntactic signals) performs significantly better then a fastText+syntax based system (a 4-6 point gap). Interestingly, supplying syntactic signals to an ELMo based system yields a small, possibly insignificant gain in performance. 


\begin{table}\label{tab:embdsyn}
\setlength{\tabcolsep}{10pt} % Default value: 6pt
\renewcommand{\arraystretch}{1.5}
\newcommand{\score}[2]{#1{ ($\pm$#2)}}
\newcommand{\acc}[1]{\multicolumn{3}{c|}{\textcolor{gray}{\rule[2pt]{0.5in}{0.5pt}} #1 \textcolor{gray}{\rule[2pt]{0.5in}{0.5pt}}}}
\newcommand{\accl}[1]{\multicolumn{3}{c}{\textcolor{gray}{\rule[2pt]{0.5in}{0.5pt}} #1 \textcolor{gray}{\rule[2pt]{0.5in}{0.5pt}}}}
\centering
\begin{tabular}{|c|cc|cc|}

\hline
\multirow{2}{*}{\backslashbox{Embeddings}{Syntax}}   & \multicolumn{2}{|c|}{\textbf{w/o Syntax}} & \multicolumn{2}{|c|}{\textbf{w/ Syntax}}  \\
& Role & Func. & Role & Func. \\ \hline
\textbf{fastText} &  \score{ 41.1 }{ 0.5 }     &  \score{ 54.0 }{ 0.5 } &  \score{ 70.9 }{ 1.6 }     &  \score{ 82.5 }{ 0.8 } \\ \hline
\textbf{ELMo} &  \score{ 78.0 }{ 1.7 }     &  \score{ 86.4 }{ 1.3 } &  \score{ 79.5 }{ 0.6 }     &  \score{ 86.8 }{ 0.9 } \\ \hline
\end{tabular}
\caption{Average role and func. prediction accuracy scores for the 4 models variations corresponding to the selection of word embeddings and use of syntax, in the gold-id/gold-syntax setting.}

\end{table}

\subsection{Errors Across the SNACS Hierarchy}

The PSS label hierarchy can be used to compute a more fine-grained measurement of disambiguation performance.  Specifically, the accuracy of predicted labels is evaluated when they are coarsened post hoc by moving up the hierarchy to a specific depth, with depth-1 representing the coarsening of the labels into the 3 root labels. The accuracy for each depth in the gold-syntax/gold-id setting is reported in table \ref{tab:coarsening}. These
results show that the classifiers often mistake a label for another that is nearby in the hierarchy.


\begin{table}[]
    \setlength{\tabcolsep}{10pt} % Default value: 6pt
    \renewcommand{\arraystretch}{1.5}
    \centering
    \begin{tabular}{|c|c|c|c|}
        \hline
            & Labels & Role & Func.\\ \hline
         Depth-4 (Exact) & 47 & 79.8 & 86.0\\ \hline
         Depth-3 & 43 & 80.2 & 86.3 \\ \hline
         Depth-2 & 26 & 83.3 & 89.8 \\ \hline
         Depth-1 & 3 & 89.6 & 94.2 \\ \hline
          
    \end{tabular}
    \caption{Accuracy scores (single model, gold identification and syntax) on the test set (480 tokens) with different levels of hierarchy coarsening of its output. “Labels” refers to the number of labels in the training set after coarsening.}
    \label{tab:coarsening}
\end{table}

\section{Discussion?}

\section{Future Work}


% Chapter Template

\chapter{Preposition Supersenses for  Prepositional Phrase Attachment} % Main chapter title

\label{Chapter3} % Change X to a consecutive number; for referencing this chapter elsewhere, use \ref{ChapterX}


%-----------------------------------
%	SECTION 1
%-----------------------------------
\section{Overview}

\section{PPA work by Belinkov et. al 2014}
\subsubsection{*Overview, relevance to this work}
This section introduces the PP-Attachment work by (Belinkov et al. 2014), which we use in our research. We incorporate PSS features onto their top-performing model, measuring the impact on accuracy. 

\subsubsection{*Intro to Belinkov work}
(Belinkov et al. 2014 todo) is perhaps the first to present a set of neural PP-Attachment models that make use of pretrained word embeddings. Their top performing model achieved state-of-the-art results, surpassing any previous attempts at this task. They experiment with different sets of word embeddings (skip-gram based and dep-gram based todo: ref), using a PP-Attachment dataset they constructed using the Penn Treebank for English and the CATiB dependency treebank (Habash and Roth, 2009) for Arabic.    

\subsubsection{* -- Isolated PPA models}
\subsubsection{* -- Incorporated to an existing parser}
The models presented operate in an isolated manner - solving PP-Attachment alone. Instead of the typical V-NP-P-NP2 formulation, which as explained in section todo was shown to be problematic, they took the more realistic approach of allowing multiple head candidates. They do, however, compare against full-scale parsers, and integrate their model into an existing one - the RGB parser (todo: ref) - resulting with an improvement of parsing accuracy in both English and Arabic when compared to using the unmodified RGB parser.

\subsubsection{* -- Model architecture in brief}
Their model architecture consists of a parameterized neural scoring function, which assigns a score to a given attachment candidate, given the corresponding word embeddings of the preposition, its child (the object) the head candidate, and their corresponding part-of-speech and WordNet, FrameNet attributes. The parameters of the scoring function are learned during training, and at inference the highest scoring candidate is chosen. They experiment with a few different scoring functions, each consists of a slightly different neural network architecture. The variations differ mostly by: (1)  the input - as one variation does not take the preposition as an input, only the head candidate and the object of the preposition; (2) the number of network layers - a single layer or two layers, and (3) the parameterization - for instance, one variation uses a different set of parameters based on the distance between the head candidate and the preposition. They also experiment with two kinds of pre-trained word embeddings: one based on the Skip-Gram model (todo: ref) and a dependency based model following (todo: ref), reporting significant accuracy boost when using the latter. Allowing the word embeddings to update during training was also reported to increase disambiguation accuracy significantly. 

\subsection{Datasets}
\subsubsection{*Structure, size and other notable stats}
\subsubsection{*Extraction process}

\subsection{Model}
\subsubsection{*Model framework}
The following section assumes that sentence words are represented using a column vector of dimension $n$. Section todo describes how such representation is constructed. 

As described in section todo, the model consists of a scoring function that assigns a score to each attachment candidate. For inference, the highest scoring candidate is selected. The detection of a preposition and its attachment candidates is done in a rule-based process as a preprocessing step, which is used for building the datasets used for training and evaluation (see section todo).

The scoring function $s(c, z, h, \theta)$ takes a preposition $z$ with child $c$, and an attachment candidate $h$ ($\theta$ is the set of learned model parameters). The score is defined to be a linear function on a vector composed from $c, z, h$:

$$s(c,z,h,\theta) = w * compose\_p(c,z,h, \theta’)$$

Where $w \in \mathbb{R}^n$ is a parameter, and $\theta'$ is the set of all parameters, excluding $w$. $compose\_p$ transforms $c, z, h$ into a single column vector $p \in \mathbb{R}^n$ using parameters in $\theta’$.  The score is obtained by taking the dot product of $w$ and $p$. 

\subsubsection{*Score function variations}
Several variations of $compose\_p$ are suggested, each corresponding to a different model variant. The top performing variant - Head-Prep-Child-Dist (HPCD) - is the one used in this research, and the only one described in detail. The rest of the variants are mainly simpler versions of HPCD.

\subsubsection{The HPCD scoring function}
The HPCD version of $compose\_p$ is defined as follows:

   $$ dist = min(word\_distance(head, prep), MAX\_DIST) $$
   $$ p_1 = tanh(W^1_{dist}[z;c] + b_{dist})   $$
   $$ p_2 = tanh(W^2_{dist}[h;p_1] + b_{dist}) $$

Where $h, z, c$ are vectors for the head, the preposition, and its child respectively; $p_1, p_2$ are the output vectors for the first and second layers respectively;
$word\_distance(head, prep)$ is the number of words between the preposition and the head candidate; and $ W^1_{dist}, W^2_{dist}, b_{dist}$ are distance-specific model parameters.

Namely, the HPCD scoring function is composed of a two-layered feed-forward neural network. The first layer takes the vectors for the preposition and its child, and the second layer takes the output vector from the first layer and the vector for the head candidate. The result is the output of the second layer. The HPCD variant takes into account the distance between the head candidate and the preposition (i.e the number of words between them), by using different neural network parameters for each distance, up to a predefined threshold (MAX\_DIST = 5). For any distance larger than the threshold, the parameters corresponding to threshold value are used.

\subsection{Word Vector Representation}

\begin{figure}
    \centering
    \includegraphics{Figures/Electron.pdf}
    \caption{An illustration of the extended word vectors for the head, preposition and child. In addition to the pretrained word embeddings, the vectors for the child and head are extended with WordNet, VerbNet and syntactic information}
    \label{fig:hpcdvectors}
\end{figure}

The word vectors for the head, preposition, and child words ($h, z, c$ in equation todo ref) are constructed using the pretrained word embeddings of the corresponding words, and extended using auxiliary syntactic information and structured knowledge resources such as WordNet and VerbNet. They experiment with two kinds of SkipGram based word embedding models. The first is the standard SkipGram model, where the context of a word is defined by the nearby words in the sentence. The second is a variation of the standard SkipGram by Bansal et al. (2014), where the context of a word is defined using its head, child and corresponding dependency labels according to the dependency tree of the given sentence. These word embeddings (referred to as \textit{syntactic} embeddings) are expected to be better at capturing dependency related information and therefore be relevant for PP-Attachment disambiguation. As their experiments show, the syntactic vectors indeed yield a significant boost in accuracy. The same 100-dimensional pretrained syntactic vectors are also used in this work. 

Both child and head vectors are extended with WordNet information. Specifically, a binary indicator is added for each top hypernym (the top the set of 26 hypernyms in top three layers of the hypernym tree). 

The head vector is also extended with binary indicators corresponding to whether the head is a noun, and in the case of a verb, whether the preposition appears in the corresponding VerbNet frame. The POS tag of the word following the head is embedded into the head vector using one-hot vector encoding. The resulting outline of the full vectors for the head, preposition and child are illustrated in figure \ref{fig:hpcdvectors}

\subsection{Training}
\subsection{Results}
\subsection{Conclusion}
\pagebreak

%-----------------------------------
%	SECTION 2
%-----------------------------------
\section{Datasets}
\subsection{World Stree Journal based dataset}
\subsubsection{*PSS enrichment of WSJ}
\subsection{STREUSLE based dataset}
\subsubsection{*Extracting PPA cases from STREUSLE}

\section{Models}
\subsection{Baseline Model: HPCD}
\subsection{HPCD extended with PSS features}

\section{Training}
\subsubsection{*Loss function and trainer remain unchanged}
\subsubsection{*Process of training on both datasets}
\subsubsection{*Hyperparameter tuning}

\section{Results}
\subsubsection{*Compare HPCD and extended HPCD on each setting and dataset}

\section{Analysis}
\subsection{Error Analysis}
\subsubsection{*Errors shared between both models}
\subsubsection{*Errors due to wrong PSS? (sampling)}

\section{Conclusions}
 
% Chapter Template

\chapter{Preposition Supersense Disambןguation in a Multilingual Setting } % Main chapter title

\label{Chapter4} % Change X to a consecutive number; for referencing this chapter elsewhere, use \ref{ChapterX}

\section{Overview}

In this chapter, we examine the transferability of the knowledge acquired by a PSS disambiguator between different languages and domains. Specifically, we evaluate the PSS model presented in chapter \ref{Chapter2} - trained solely using the STREUSLE corpus of English web reviews - on the recently introduced Mandarine Chinese PSS corpus by \cite{chinesecorpus}. The Chinese corpus consists of the first 20 chapters of \textit{The Little Prince}, along with manually annotated PSS labels for all identified adpositions. First, words in the Chinese corpus are automatically aligned to the corresponding words in the parallel English translation (section \ref{sec:chinesecorpus}). As the word alignment is noisy and incomplete, the PSS model and training process are slightly altered to better handle the information gaps (section \ref{sec:pssmodelmod}). The PSS disambiguation accuracy on the Chinese corpus is reported in section \ref{sec:zhpssres}. The near-baseline levels of disambiguation accuracy on the Chinese corpus confirm the significance of domain and language for model performance.

\section{Aligning the Parallel Corpus} \label{sec:chinesecorpus}

The Chinese corpus was word-aligned to the corresponding English translation using TsinghuaAligner (\cite{liu2014contrastive}) - an unsupervised log-linear model trained using a \textit{Contrastive Learning} objective: assigning higher probabilities to correct sentence pairs, than to noisy ones generated by shuffling pairs in the training set. 

The alignment process is rather noisy and incomplete - out of 518 annotated Chinese prepositions, only 98 (19\%) were aligned to English prepositions. 49 (9\%) Chinese prepositions were aligned to non-prepositional English words, and the remaining 371 (72\%) prepositions did not receive an alignment at all. The PSS disambiguator also considers the governor and object of each preposition, but since these are not available for Chinese, it was modified to use the parent and grandparent of the preposition according to the UD tree, which provide a close approximation. The parent and grandparent were determined according to the English UD tree in the case of aligned prepositions, and the Chinese UD tree in the case of unaligned prepositions. The UD trees were obtained using the CoreNLP parsers for Chinese and English (\cite{manning14stanford}). 

The parents and grandparents of an unaligned preposition, were also not fully aligned: for the 371 unaligned Chinese prepositions, 138 (37\%) of the parents and 162 (43\%) of the grandparents remained unaligned. The remaining auxiliary information required by the PSS model (UD labels, POS, NER, etc.) was obtained by applying the CoreNLP parser on the English corpus. 

\section{PSS Model and Training Adjustments} \label{sec:pssmodelmod}

The PSS model presented in Chapter \ref{Chapter2} was slightly altered to better handle the noisy alignment of words between the Chinese and the English corpora. As previously mentioned, the model uses the parent and grandparent of a preposition instead of the governor and object. For the model to handle unaligned prepositions, they are randomly replaced with "blank" tokens during training, in a tunable probability. "blank" tokens are not fed into the LSTM layer, and thus do not have a contextualized representation. In such cases, a single (trainable) parameter vector is used as the contextualized representation portion of the enriched preposition vector (section \ref{sec:enrichedvec}). Similarly, the parent and grandparent contextualized representations portion of the enriched preposition vector are randomly dropped, by replacing them with a single parameter vector, representing a missing parent/grandparent. 

\section{Results} \label{sec:zhpssres}

\begin{table}[]
    % \setlength{\tabcolsep}{6pt} % Default value: 6pt
    \newcommand\Tstrut{\rule{0pt}{2.6ex}}       % "top" strut
    \newcommand\Bstrut{\rule[-0.9ex]{0pt}{0pt}} % "bottom" strut
    \newcommand{\TBstrut}{\Tstrut\Bstrut} % top&bottom struts

    \renewcommand{\arraystretch}{1}
    \newcommand{\score}[2]{#1 {\footnotesize ($\pm$#2)}}
    \newcolumntype{C}{c}
    \centering
    \begin{tabular}{*{9}{C}}
         \toprule
         System    & \multicolumn{2}{C}{STREUSLE acc.} & \multicolumn{6}{C}{Chinese acc.}  \\ \cmidrule(lr){4-9}
                   &    \multicolumn{2}{C}{} & \multicolumn{2}{C}{All (518)}      & \multicolumn{2}{C}{Aligned (147)} & \multicolumn{2}{C}{\makecell{Aligned To\\Prep. (98)}} \\
         \cmidrule(lr){2-3}
         \cmidrule(lr){4-5}
         \cmidrule(lr){6-7}
         \cmidrule(lr){8-9}
                                    & Role & Fxn.     & Role  & Fxn.   & Role  & Fxn.  & Role  &  Fxn. \Bstrut \\
         \hline

         MF-Prep & 40.62 & 53.33 & 36.67 & 41.69 & 48.29 & 49.65 & 68.36 & 69.38  \Tstrut \\
         Neural  & 79.69 & 88.74 & 36.1 & 36.29 & 52.38 & 51.02 & 68.36 & 66.32  \\
         Neural + MF-Prep  & - & - & 37.83 & 42.08 & - & - & - & - \\
         \bottomrule
    \end{tabular}
    \caption{PSS disambiguation accuracy of the MF-Prep baseline, the neural model, and a combined model, on the STREUSLE test set, and the Chinese corpus. For Chinese, we also report the accuracy over the subset of prepositions that received and alignment, and the subset which was aligned to English prepositions.}
    \label{tab:chineseresults}
\end{table}

The modified PSS model was trained as described in Chapter \ref{Chapter2}, using the STREUSLE training set for training, and the development set for hyperparameter tuning. The model, trained solely on English sentences, is then evaluated on the Chinese corpus with the use of the word-aligend, parallel English translation. The disambiguation accuracy of the model, as well as the MF-Prep baseline (section \ref{sec:pssbaselines}), are reported in table \ref{tab:chineseresults}. We also report the overall Chinese accuracy of a combined model, which uses the neural model for aligned prepositions, and the MF-Prep baseline model for the rest.

As expected, the disambiguation accuracy is highly dependent on the quality of the alignment: Chinese prepositions aligned to English prepositions are almost twice as accurate than the average Chinese preposition. Still, the change of domain and language have a critical negative impact on the neural model performance (a 40-50 pt. drop overall, and a 10-20 pt. drop on the subset of Chinese prepositions aligned to English prepositions). Moreover, the neural model accuracy is of the same order of magnitude as the MF-Prep baseline across all reported Chinese subsets - indicating that any complex pattern it might have detected in the training data was not as useful for Chinese predictions. The combined model (Neural + MF-Prep) performs similarly to both models, which is another indication that they behave similarly on Chinese.










 
\chapter{Conclusions}

*1. PSS labels can be predicted with a high accuracy using a neural model
*2. PSS signal may be useful for PP Attachment
*3. PSS may be useful for normalizing different domains and languages
*4. (a) future work - further assess PSS contribution to pp-att
(b) and to other preposition related tasks: translation, entity-relation extraction
(c) integrate with additional schemes both to facilitate a more complete scheme, and for join-learning.
(d) further experiment with the cross lingual setting - better alignment, more data, multilingual embeddings 

%----------------------------------------------------------------------------------------
%	THESIS CONTENT - APPENDICES
%----------------------------------------------------------------------------------------

\appendix % Cue to tell LaTeX that the following "chapters" are Appendices

% Include the appendices of the thesis as separate files from the Appendices folder
% Uncomment the lines as you write the Appendices

% Appendix A

\chapter{Hyperparameters for the PSS Neural Model (Chapter 2)} % Main appendix title

\label{AppendixA} % For referencing this appendix elsewhere, use \ref{AppendixA}

Table \ref{tab:hyperparamsA} lists the hyperparameters used by the
PSS neural system presented in chapter \ref{Chapter2}, tuned separately for each of the four settings corresponding to gold/automatic auxiliary syntax and gold/automatic target identification.

\begin{table}[]
  \small
	\centering
	\begin{tabular}{@{}l|cccc@{}}
		\toprule
		Hyperparam. & Auto ID/Auto Parse & Auto ID/Gold Parse & Gold ID/Auto Parse & Gold ID/Gold Parse\\
		\midrule
        Embedding Method & ELMo & ELMo & ELMo & ELMo \\
        ELMo layer & 1 & 1 & 1 & 1 \\
        Token internal embd. dimension & 0 & 0 & 10 & 0 \\
        Use Gov-Obj? & No & Yes & Yes & Yes \\
        POS embd. dimension  & 0 & 0 & 0 & 10 \\
        UD dependencies embd. dimension  & 0 & 5 & 0 & 0 \\
        NER  embd. dimension  & 0 & 5 & 0 & 5 \\
        GOVOBJ-CONFIG embd. dimension  & 0 & 3 & 3 & 3 \\
        LEXCAT embd. dimension  & 3 & 3 & 3 & 3 \\
        MLP layers  & 2 & 0 & 2 & 2 \\
        MLP layer dimension  & 100 & - & 200 & 200 \\
        MLP activation  & tanh & - & relu & relu \\
        MLP Dropout Prob.  & 0.46 & - & 0.14 & 0.34 \\
        LSTM layers  & 1 & 2 & 1 & 2 \\
        LSTM hidden layer dimension  & 200 & 200 & 200 & 200 \\
        Is BiLSTM? & Yes & Yes & Yes & Yes \\
        LSTM Dropout Prob.  & 0.49 & 0.43 & 0.42 & 0.4 \\
        Learning rate  & 0.15 & 0.15 & 0.2 & $10^{-1}$ \\
        Learning rate decay  & 0 & 0 & $10^{-4}$ & 0 \\
		\bottomrule
	\end{tabular}
	\caption{\label{tab:hyperparamsA}
        Selected hyperparameters of the neural PSS disambiguation system for each of the four settings. The parameters were tuned over the STREUSLE development set using random grid search.
	}

\end{table}

% Appendix B

\chapter{Hyperparameters for HPCD variants (Chapter 3)} % Main appendix title

\label{AppendixB} % For referencing this appendix elsewhere, use \ref{AppendixA}

Table \ref{tab:hyperparamsB} lists the hyperparameters used by the PP-Attachment models presented in chapter \ref{Chapter3} - vanilla HPCD, and its PSS-bearing variant (HPCD-PSS). The hyperparameters were tuned separately in each of the four settings corresponding to the selection of training set (PTB/PTB+STREUSLE) and model (HPCD/HPCD-PSS). 

\begin{table}[]
  \small
	\centering
	\begin{tabular}{@{}l|cccc@{}}
		\toprule
		Hyperparam. & HPCD/PTB & HPCD-PSS/PTB & HPCD/PTB+STREUSLE & HPCD-PSS/PTB+STREUSLE\\
		\midrule
            $p_1$ dimension & 100 & 100 & 100 & 100 \\
%            & 1 & 1 & 1 & 1 \\
            $p_2$ dimension & 100 & 100 & 100 & 100 \\
%            & 1 & 1 & 1 & 1 \\
            Activation & tanh & tanh & tanh & tanh \\
            Use PSS? & No & Yes & No & Yes \\
            PSS Embd. type & - & Lookup Table & - & Lookup Table \\
            PSS Embd. dimension& - & 50 & - & 50 \\
            Dropout Prob. & 0.5 & 0.5 & 0.5 & 0.5 \\
            Learning rate & 0.1 & 0.1 & 0.1 & 0.1 \\
            Learning rate decay& 0 & 0 & 0 & 0 \\
            Update embeddings? & Yes & Yes & Yes & Yes \\
%            & No & Yes & No & Yes \\
%            & SimpleSGDTrainer & SimpleSGDTrainer & SimpleSGDTrainer & SimpleSGDTrainer \\
%            & 100 & 100 & 100 & 100 \\
		\bottomrule
	\end{tabular}
	\caption{\label{tab:hyperparamsB}
        Selected hyperparameters of the PP-Attachment systems for each of the four settings correspoding to the training set (PTB/PTB+STREUSLE) and model (HPCD/HPCD-PSS).
	}

\end{table}


% Appendix C

\chapter{Hyperparameters for the PSS Neural Model (Chapter 4)} % Main appendix title

\label{AppendixC} % For referencing this appendix elsewhere, use \ref{AppendixA}

% \begin{table}[h]
%   \small
% 	\centering
% 	\begin{tabular}{@{}l|cccc@{}}
% 		\toprule
% 		Hyperparam. & \thead{Auto ID\\Auto Parse} & \thead{Auto ID\\Gold Parse} & \thead{Gold ID\\Auto Parse} & \thead{Gold ID\\Gold Parse}\\
% 		\midrule
%         Embedding Method & ELMo & ELMo & ELMo & ELMo \\
%         ELMo layer & 1 & 1 & 1 & 1 \\
%         Token internal embd. dimension & 0 & 0 & 10 & 0 \\
%         Use Gov-Obj? & No & Yes & Yes & Yes \\
%         POS embd. dimension  & 0 & 0 & 0 & 10 \\
%         UD dependencies embd. dimension  & 0 & 5 & 0 & 0 \\
%         NER  embd. dimension  & 0 & 5 & 0 & 5 \\
%         GOVOBJ-CONFIG embd. dimension  & 0 & 3 & 3 & 3 \\
%         LEXCAT embd. dimension  & 3 & 3 & 3 & 3 \\
%         MLP layers  & 2 & 0 & 2 & 2 \\
%         MLP layer dimension  & 100 & - & 200 & 200 \\
%         MLP activation  & tanh & - & relu & relu \\
%         MLP Dropout Prob.  & 0.46 & - & 0.14 & 0.34 \\
%         LSTM layers  & 1 & 2 & 1 & 2 \\
%         LSTM hidden layer dimension  & 200 & 200 & 200 & 200 \\
%         Is BiLSTM? & Yes & Yes & Yes & Yes \\
%         LSTM Dropout Prob.  & 0.49 & 0.43 & 0.42 & 0.4 \\
%         Learning rate  & 0.15 & 0.15 & 0.2 & $10^{-1}$ \\
%         Learning rate decay  & 0 & 0 & $10^{-4}$ & 0 \\
% 		\bottomrule
% 	\end{tabular}
% 	\caption{\label{tab:hyperparamsC}
%         Selected hyperparameters of the neural PSS disambiguation system presented in chapter \ref{Chapter4}.
% 	}

% \end{table}

\begin{table}[h]
  \small
	\centering
	\begin{tabular}{@{}l|c@{}}
		\toprule
		Hyperparam. & Value \\
		\midrule
        Embedding Method & ELMo \\
        ELMo layer & 1 \\
        Token internal embd. dimension & 0 \\
        Use Parent? & Yes \\
        Use Grandparent? & Yes \\
        POS embd. dimension  & 10 \\
        UD dependencies embd. dimension  & 0 \\
        NER  embd. dimension  & 0 \\
        MLP layers  & 0 \\
        MLP layer dimension  & - \\
        MLP activation  & - \\
        MLP Dropout Prob.  & - \\
        Parent Dropout Prob. & 0.2 \\
        Grandparent Dropout Prob. & 0.3 \\
        Prep. Dropout Prob. & 0.3 \\
        LSTM layers  & 0 \\
        LSTM hidden layer dimension  & - \\
        Is BiLSTM? & - \\
        LSTM Dropout Prob.  & - \\
        Learning rate  & 0.15 \\
        Learning rate decay  & 0 \\
		\bottomrule
	\end{tabular}
	\caption{\label{tab:hyperparams}
        Selected hyperparameters of the neural system, tuned on the STREUSLE development set.
	}

\end{table}

Table \ref{tab:hyperparamsC} lists the hyperparameters used by the
PSS neural system presented in chapter \ref{Chapter4}, trained and tuned on the STREUSLE PSS corpus in English, and evaluated on the Chinese PSS corpus based on \textit{The Little Prince}.



%----------------------------------------------------------------------------------------
%	BIBLIOGRAPHY
%----------------------------------------------------------------------------------------

\printbibliography[heading=bibintoc]

%----------------------------------------------------------------------------------------

\end{document}  
